\subsection{适度函数与适度测度}

命题 5. --- 设 A 是 T 的子集；下列性质等价：

\begin{enumerate}
\def\labelenumi{\alph{enumi})}
\item
  集 A 含于一个 \(\mu\) -可积开集序列的并中。
\item
  集 A 含于一个 \(\mu\) -可积集序列的并中。
\item
  集 A 含于一个紧集序列与一个 \(\mu\) -可忽略集的并中。
\end{enumerate}

显然性质 a) 与 c) 各自蕴含 b)。反之，b) 蕴含 a)，因为每个外测度有限的集都含于一个可积开集（第四章，§1 第 4 号命题 19）；b) 蕴含 c)，因为每个可积集是一个紧集序列与一个可忽略集之并（第四章，§4 第 6 号定理 4 推论 2）。

定义 2. --- T 的子集若满足命题 5 的等价条件，就称为 \(\mu\) -适度的。定义在 T 上、取值于某个向量空间或 \(\overline{R}\) 中的函数，若它在 T 的一个 \(\mu\) -适度子集的余集上为零，就称为 \(\mu\) -适度的。若 T 是 \(\mu\) -适度集，则称测度 \(\mu\) 是适度的。

若 \(\mu\) 是适度测度，则 T 上每个函数都是 \(\mu\) -适度的，且 T 的每个子集都是 \(\mu\) -适度的。

注记. --- 1) 若 \(\theta\) 是 T 上的复测度，则称函数 \(f\) 是 \(\theta\) -适度的（相应地，称 \(\theta\) 是适度的），如果 \(f\) 是 \(|\theta|\) -适度的（相应地，如果 \(|\theta|\) 是适度的）。

\begin{enumerate}
\def\labelenumi{\arabic{enumi})}
\setcounter{enumi}{1}
\item
  每个有界测度都是适度的；若 \(\mathrm{T}\) 是紧集的可数并，则 \(\mathrm{T}\) 上每个测度都是适度的。
\item
  设 \((f_n)\) 是 \(\mu\) -适度函数的序列，取值于 \(\overline{\mathbf{R}}\) 。对每个 \(n\) ，设 \(\mathrm{U}_n\) 是一个开集，它是外测度有限的开集的可数并，使 \(f_n\) 在 \(\mathrm{U}_n\) 外为零。于是函数 \(s = \sum_{n \in \mathbf{N}} |f_n|\) 在 \(\bigcup_{n\in \mathbf{N}}\mathrm{U}_n\) 外为零；因而它是 \(\mu\) -适度的，同样的结论对
\end{enumerate}

所有函数 \(f\) （满足 \(|f| \leqslant s\) ）成立。这特别适用于函数 \(\liminf_{n \to \infty} f_n\) 、\(\limsup_{n \to \infty} f_n\) 与 \(\sum_{n \in \mathbf{N}} f_n\) （若和有定义）。

\begin{enumerate}
\def\labelenumi{\arabic{enumi})}
\setcounter{enumi}{3}
\tightlist
\item
  几乎处处等于一个适度函数的函数是适度的。
\end{enumerate}

命题 6. --- 设 f 是定义在 T 上的正数值函数，且是 \(\mu\) -可测的与 \(\mu\) -适度的。则存在一个序列 \((h_{n})_{n\in\mathbf{N}}\) （由 \(\mathcal{F}_{+}(\mathrm{T})\) 的元素组成），其和等于 f，并具有下列性质：

\begin{enumerate}
\def\labelenumi{\arabic{enumi})}
\item
  函数 \(h_0\) 是 \(\mu\) -可忽略的。
\item
  对每个 \(n \geqslant 1\) ，存在紧集 \(K_{n}\) ，使 \(h_{n}\) 在 \(K_{n}\) 外为零，且 \(h_{n}\) 在 \(K_{n}\) 上的限制是有限且连续的。
\end{enumerate}

假设 \(f\) 是正可测函数的序列 \((f_n)\) 之和，其中每个函数都具有命题中所述的性质；显然 \(f\) 也具有该性质。令

\[
f _ {n} = \inf (f, n + 1) - \inf (f, n)
\]

对每个 \(n \in \mathbf{N}\) ；由于 \(f\) 等于序列 \((f_n)\) 之和，故只需在假定 \(f\) 适度且有界的情形下证明该命题。于是用 \(A\) 表示由 \(t \in T\) （使 \(f(t) > 0\) 者）组成的集；\(A\) 可测且适度，因此存在两两不相交的可积集的序列 \((A_n)\) 使 \(A = \bigcup_{n} A_n\) 。我们可化归

为对函数 \(f\varphi_{\mathrm{A}_n}\) 证明命题；换言之，可设 \(f\) 有界且在一个可积集 I 外为零。但 I 是一个可忽略集 N 与两两不相交的紧集的序列 \((\mathbf{L}_n)\) 之并（第四章，§4 第 6 号定理 4 推论 2）。于是化归为处理 \(f\) 有界且在一个紧集 L 外为零的情形。

设 \(\mathfrak{K}\) 为紧子集 \(K\) （\(T\) 中使 \(f|K\) 连续者）所成之集；由于 \(\mathfrak{K}\) 是 \(\mu\) -稠密的（第四章，§5 第 10 号命题 15），\(L\) 是一个可忽略集 \(N\) 与两两不相交元素的一个序列 \((K_n)_{n \geqslant 1}\) （\(\mathfrak{K}\) 中的）之并（第四章，§5 第 8 号定义 6）。于是函数 \(h_0 = f\varphi_N\) 与 \(h_n = f\varphi_{K_n}\) （对 \(n \geqslant 1\) ）满足命题中的条件。

下述命题使我们可以把上积分的研究化归为本质上积分的研究。

命题 7. --- 设 \(f\) 是 \(\mathcal{F}_{+}(\mathrm{T})\) 的一个元素。

\begin{enumerate}
\def\labelenumi{\arabic{enumi})}
\item
  若函数 \(f\) 不是 \(\mu\) -适度的，则 \(\mu^{*}(f) = +\infty\) 。
\item
  若函数 \(f\) 是 \(\mu\) -适度的，则 \(\mu^{*}(f) = \mu^{\bullet}(f)\) 。
\item
  若 \(\mu^{\bullet}(f) < +\infty\) ，则存在一个 \(\mu\) -适度子集 A，它是 T 的紧子集的一个序列之并，使 \(f = f\varphi_{\mathrm{A}}\) 局部几乎处处成立。
\end{enumerate}

第一个断言由第四章，§5 第 6 号引理 1 立即得出。为建立第二个断言，用 A 表示一个适度子集，使 f 在 A 外为零；A 是一个可忽略集 \(A_{0}\) 与紧集的序列 \((\mathrm{A}_{n})_{n\geqslant1}\) （可设它是递增的）之并。于是函数 f 几乎处处等于诸函数 \(f\varphi_{\mathrm{A}_{n}}\) （ \(n\geqslant1\) ）的上包络，因此（第四章，§1 第 3 号定理 3 及 §2 第 3 号命题 6）

\[
\mu^ {*} (f) = \lim _ {n \rightarrow \infty} \mu^ {*} (f \varphi_ {\mathrm{A} _ {n}}) \leqslant \mu^ {\bullet} (f),
\]

于是由公式 (1) 得等式 \(\mu^{*}(f)=\mu^{\bullet}(f)\) 。最后，设 \(\mu^{\bullet}(f)<+\infty\) ；存在紧集的递增序列 \((\mathrm{A}_{n})\) 使

\[
\mu^ {\bullet} (f) = \sup _ {n} \mu^ {*} (f \varphi_ {\mathrm{A} _ {n}}).
\]

令 \(A = \bigcup_{n} A_{n}\) ；右端等于 \(\mu^{*}(f\varphi_{\mathrm{A}})\) （第四章，§1 第 3 号定理 3），即等于 \(\mu^{\bullet}(f\varphi_{\mathrm{A}})\) （由命题 1，或由上面的 2)）。由于 \(\mu^{\bullet}(f) = \mu^{\bullet}(f\varphi_{\mathrm{A}}) + \mu^{\bullet}(f\varphi_{\mathbb{C}_{\mathrm{A}}})\) （命题 2），我们有 \(\mu^{\bullet}(f\varphi_{\mathbb{C}_{\mathrm{A}}}) = 0\) ，由此得出 3)。

推论 1. --- 为使 f 可忽略，必须且只须它局部可忽略且适度。

推论 2. --- 若 \(\mu\) 是适度测度（特别地，若 \(\mu\) 有界，或 T 在无穷远处可数），则 \(\mu^{*} = \mu^{\bullet}\) 。

命题 8. --- a) 设 H 是由 ≥ 0 的下半连续函数组成的集，对关系 ≤ 有向；则

\[
\mu^ {\bullet} \left(\sup _ {h \in \mathrm{H}} h\right) = \sup _ {h \in \mathrm{H}} \mu^ {\bullet} (h).
\]

\begin{enumerate}
\def\labelenumi{\alph{enumi})}
\setcounter{enumi}{1}
\tightlist
\item
  设 H 是由 \(\geqslant0\) 的上半连续函数组成的集，对关系 \(\geqslant\) 有向；若 H 中存在函数 \(h_{0}\) 使 \(\mu^{\bullet}(h_{0})<+\infty\) ，则
\end{enumerate}

\[
\mu^ {\bullet} \Bigl (\inf _ {h \in \mathrm{H}} h \Bigr) = \inf _ {h \in \mathrm{H}} \mu^ {\bullet} (h).
\]

鉴于命题 4，断言 a) 是第四章，§1 第 1 号定理 1 的重复。为建立 b)，令 \(\eta = \inf_{h\in \mathrm{H}}h\) ，并设 \(a\) 是一个数 \(>0\) 。存在紧集 K 使（第四章，§4 第 4 号命题 5 推论 1）：

\[
\mu^ {\bullet} (h _ {0}) - a \leqslant \mu^ {*} (h _ {0} \varphi_ {\mathrm{K}}) = \mu (h _ {0} \varphi_ {\mathrm{K}}) \leqslant \mu^ {\bullet} (h _ {0}).
\]

函数 \(h\varphi_{K}\) （其中 h 取遍 H）组成一个上半连续函数的集，对关系 \(\geqslant\) 有向，且包含一个可积函数。因此（第四章，§4 第 4 号命题 5 推论 2）：

\[
\mu^ {*} (\eta \varphi_ {\mathrm{K}}) = \inf _ {h \in \mathrm{H}} \mu^ {*} (h \varphi_ {\mathrm{K}}).
\]

但是（第四章，§4 第 4 号命题 5 推论 1）\(\mu^{\bullet}(h_{0}\varphi_{\mathbb{C}_{\mathbf{K}}}) \leqslant a\) ，从而 \(\mu^{\bullet}(h\varphi_{\mathbb{C}_{\mathbf{K}}}) \leqslant a\) 对每个函数 \(h \in \mathrm{H}\) （满足 \(h \leqslant h_{0}\) ）成立。因此最后有：

\[
\mu^ {\bullet} (\eta) \geqslant \mu^ {*} (\eta \varphi_ {\mathrm{K}}) = \inf _ {h \in \mathrm{H}} \mu^ {*} (h \varphi_ {\mathrm{K}}) \geqslant \inf _ {h \in \mathrm{H}} \mu^ {\bullet} (h) - a.
\]

由于不等式 \(\mu^{\bullet}(\eta) \leqslant \inf_{h \in \mathrm{H}} \mu^{\bullet}(h)\) 是明显的，且 a 任意，命题得证。

